\documentclass[10pt,article]{article}
\usepackage[a4paper,margin=20mm]{geometry}
\usepackage{amsmath,amssymb}
\usepackage{mathrsfs}
\usepackage{amsthm}
\usepackage[colorlinks=true]{hyperref}
\usepackage{booktabs}
\providecommand{\tightlist}{\setlength{\itemsep}{0pt}\setlength{\parskip}{0pt}}
\setcounter{secnumdepth}{-1}
\emergencystretch=3em
\allowdisplaybreaks
\newcommand{\rulebar}{\begin{center}\rule{0.5\linewidth}{0.5pt}\end{center}}

\begin{document}

\begin{center}
{\Large\bfseries Computational Viability Certification under Incomplete Observation}\\[0.4em]
{\large\bfseries Supplementary Material: Complete Proofs and the Verification Record}\\[0.6em]
Amin Abaee\\
{\small Independent Researcher}\\
{\small September 24, 2026}
\end{center}

\section*{S1. Complete proofs}

Throughout, labels are \(a = (\theta, i, t) \in \mathscr S = \Theta \times
\{1,\dots,s\} \times [0,T]\), \(k_a\), \(\beta_a\), \(\mathcal F_a\),
\(J_{\mathcal I}(T)\), \(\phi_U\), \(\Gamma_{\mathcal I}\),
\(\rho_{\mathcal H,G}\), and the program (5) are as in the main text.

\textbf{Proof of Proposition 1 (robust rows).} The linear solution formula
gives, for \(t \le T\),
\[
c_i^\top x_\theta(t) = c_i^\top X_\theta(t,0)x_0 + \int_0^t c_i^\top
X_\theta(t,\sigma) g_\theta(\sigma)\,d\sigma + \int_0^t k_a(\sigma)^\top
u(\sigma)\, d\sigma + \int_0^t c_i^\top X_\theta(t,\sigma)E_\theta(\sigma)
d(\sigma)\, d\sigma .
\]
The initial-state term is bounded over \(x_0 \in P_\theta\) by the support
value \(h_{P_\theta}(X_\theta(t,0)^\top c_i)\), and the disturbance term,
maximized over all measurable \(d(\cdot)\) with values in \(D_\theta\), is
\(\int_0^t h_{D_\theta}(E_\theta(\sigma)^\top X_\theta(t,\sigma)^\top
c_i)\, d\sigma\); for polytopic \(D_\theta\) a maximizing vertex selection is
measurable, so the bound is attained by an admissible disturbance signal.
Collecting terms, the constraint \(c_i^\top x_\theta(t) \le b_i\) holds for
every \((x_0, d)\) if and only if \(\mathcal F_a(\pi) \le 0\), and the
separation of the maximizations from the control term is exactly that the
control signal in a mode does not depend on \(x_0\) or \(d\).
\rulebar

\textbf{Proof of Proposition 2 (value and nonviability).} The admissible
signal space is a finite product of the sets \(L^\infty(I_r; U)\); because
\(U\) is compact and convex these are weak-* compact by the Banach--Alaoglu
theorem. The safety functionals are affine and weak-* continuous because
their kernels \(k_a\) lie in \(L^1\). Their supremum over the finitely many
labels is therefore weak-* lower semicontinuous and attains its minimum on
the policy space: \(J_{\mathcal I}(T)\) is well defined. Safety of some
policy is \(\max_a \mathcal F_a(\pi) \le 0\), so \(B_0 \in
\mathcal K_{\mathcal I}^T\) iff \(J_{\mathcal I}(T) \le 0\); since safety at
the margin is viability (constraints closed) and \(J_{\mathcal I}(T) = 0\)
admits safe policies, nonviability is \(J_{\mathcal I}(T) > 0\). The strict
inequality is thus a consequence of compactness and closedness, not an added
assumption.
\rulebar

\textbf{Proof of Theorem 4 (continuous-to-finite).}
\emph{Step 1 (control-moment inequality).} For a block \(J_j\) and any
measurable \(u\) with values in \(U\),
\[
\Bigl(\int_{J_j} k_a\Bigr)^{\!\top} \frac{1}{|J_j|}\int_{J_j} u
- \int_{J_j} k_a^\top u
= \int_{J_j} (\bar k_{a,\mathcal H} - k_a)^\top u
\le \int_{J_j} h_U(\bar k_{a,\mathcal H} - k_a),
\]
using \(\sup_{v \in U} w^\top v = h_U(w)\). Summing over blocks proves (3);
the averages remain compatible with the information structure because the
mesh contains all sensing times, so common-history controls use the same
block variable. For the bound, choose a center \(u_c\) of \(U\); since
\(\int_{J_j}(\bar k - k) = 0\), the center cancels and
\(e_{a,\mathcal H} \le R_U \|k_a - \bar k_{a,\mathcal H}\|_{L^1}\). A
bounded-variation function satisfies \(\|k_a - \bar k_{a,\mathcal
H}\|_{L^1} \le h_u\, \mathrm{TV}(k_a)\), whence \(e_{a,\mathcal H} \le R_U
V_U h_u\). The pointwise bound behind it: for every \(\sigma \in J\),
\(\|\bar k - k(\sigma)\| = \bigl\|\tfrac{1}{|J|}\int_J (k(\tau) -
k(\sigma))\,d\tau\bigr\| \le \mathrm{TV}(k; J)\), the norm of any
increment being at most the total variation; integrate over \(J\) and
sum. The same estimate bounds the reflected integrand
\(\int h_U(k_a - \bar k)\) (center cancellation is symmetric), which
is what Step 3 uses.

\emph{Step 2 (lower bound of (8)).} Take any admissible continuous-time
policy and form its block averages, which lie in \(U\) and satisfy (3). For
every grid label,
\[
A_a v - \beta_a^{+} - \bar e_a \le \int_0^T k_a^\top u^\pi_\theta -
\beta_a = \mathcal F_a(\pi),
\]
so the finite program attains a value no larger than the maximum robust
violation of that policy (the display uses the two-sided moment estimate
of Step 1, \(|A_a v - \int k_a^\top u^\pi| \le \bar e_a\), and
\(\beta_a^{+} \ge \beta_a\)). Minimizing over policies gives
\(\rho_{\mathcal H,G} \le J_{\mathcal I}(T)\) --- the block-moment
vector of every admissible policy is feasible for the finite program, at
objective \(\max_a \mathcal F_a(\pi)\).

\emph{Step 3 (upper bound of (8)).} Take an optimizer \(v^{*}\) of (5) and
implement the corresponding piecewise-constant information-adapted control,
for which \(\int_0^T k_a^\top u^{v^{*}}_\theta = A_a v^{*}\) --- exactly,
with no moment error, because the implemented control is constant on each
block; the implementation price is paid only through the grid labels
versus all times below. At grid times,
\(\mathcal F_a(v^{*}) \le \rho_{\mathcal H,G} + \bar e_a + \delta_\beta\);
the Lipschitz hypothesis (6) extends the bound to every time, giving
\(\max_a \mathcal F_a(v^{*}) \le \rho_{\mathcal H,G} + R_U V_U h_u +
\delta_\beta + L h_t\). Since \(J_{\mathcal I}(T)\) minimizes over all
admissible policies, the upper bound follows. Together with (4) this gives
the sharpened form. The asymmetry --- averaging arbitrary controls proves
the lower bound, implementing piecewise-constant controls proves the upper
bound --- is what makes na\"ive infeasibility-by-discretization arguments
unsound.

\emph{Step 4 (sound finite certificates).} Write the block input constraints
as \(F_{\mathrm{blk}} v \le f_{\mathrm{blk}}\). The program's dual is
\[
\max_{\lambda,\mu}\ -\lambda^\top(\beta^{+} + \bar e) - \mu^\top
f_{\mathrm{blk}}
\quad \text{s.t.} \quad
A^\top \lambda + F_{\mathrm{blk}}^\top \mu = 0,\quad
\lambda \ge 0,\ \mathbf 1^\top \lambda = 1,\ \mu \ge 0,
\tag{12}
\]
so a positive dual objective certifies infeasibility of \(\{A v \le
\beta^{+} + \bar e,\ v \in U^{Q}\}\). By Step 1 every genuinely safe
measurable controller would generate a feasible \(v\); the finite
infeasibility certificate therefore excludes every such controller. Moreover
the weighted continuous-time violation is at least the dual objective: a
normalized certificate with margin \(\gamma > 0\) guarantees an actual facet
violation of at least \(\gamma\) in one stored test (Remark 1 of the main
text).

\emph{Step 5 (completeness).} If the system is nonviable, let
\(\gamma = J_{\mathcal I}(T) > 0\) and choose meshes making the error in
(8) smaller than \(\gamma\); then \(\rho_{\mathcal H,G} > 0\) and Step 4's
finite LP duality yields a positive certificate, proving (iii). No
belief-state recursion is used.

\emph{Step 5b (exact atomic dual representation, proof of (9)).} Write
\(\Pi\) for the policy space and \(\mathscr S\) for the label set.

(a) \emph{Policy compactness.} Each factor
\(\{u \in L^\infty(I_r; U)\}\) is weak-* closed --- \(U\) is a finite
intersection of halfspaces \(f_i^\top v \le g_i\), and for
\(f_i \mathbf 1_{I_r} \in L^1\) the set \(\{u : f_i^\top u(\sigma)
\le g_i \text{ a.e.}\}\) is weak-* closed --- and weak-* bounded;
Banach--Alaoglu makes it weak-* compact and convex. The finite product
\(\Pi\) is therefore weak-* compact, convex, and metrizable (each
\(L^1\) is separable, so bounded weak-* sets are metrizable).

(b) \emph{Joint continuity of the label-policy pairing.} Each
\(\mathcal F_a\) is a finite sum of \(L^1\)-pairings against cell
signals minus a constant, so \(\pi \mapsto \mathcal F_a(\pi)\) is weak-*
continuous. Under the label-continuity hypothesis
(\(\sup_\sigma \|k_a - k_{a'}\| \to 0\), \(|\beta_a - \beta_{a'}| \to 0\)),
\[
\sup_\pi \, |\mathcal F_a(\pi) - \mathcal F_{a'}(\pi)|
\;\le\; T \, \max_{v \in U} \|v\| \, \sup_\sigma \|k_a - k_{a'}\|
\;+\; |\beta_a - \beta_{a'}| \;\to\; 0,
\]
uniformly in \(\pi\); with separate continuity this gives joint
continuity of \((a, \pi) \mapsto \mathcal F_a(\pi)\) on the compact
\(\mathscr S \times \Pi\), in particular boundedness.

(c) \emph{Mixed extension and Sion.} Let
\(\mathcal P(\mathscr S)\) be the Borel probability measures with the
weak-* topology: by the Riesz representation theorem it is compact,
convex, and metrizable. The bilinear
\(G(\pi, \mu) = \int \mathcal F_a(\pi)\,d\mu(a)\) is affine in
each variable --- hence simultaneously quasi-convex and quasi-concave ---
and continuous in each variable (\(\pi\)-continuity by dominated
convergence against the fixed finite \(\mu\); \(\mu\)-continuity by
pairing against the continuous \(a \mapsto \mathcal F_a(\pi)\)).
Sion's minimax theorem (both spaces compact convex) gives
\(\min_\pi \max_\mu G = \max_\mu \min_\pi G\) with the outer
maximum attained; and \(\max_\mu G(\pi, \cdot) = \max_a
\mathcal F_a(\pi)\), because the integral of a continuous function
over a probability measure lies within its min--max range while the Dirac
measures attain both ends. Hence
\(J_{\mathcal I}(T) = \max_\mu \min_\pi G\).

(d) \emph{Continuity of the value function.} \(g(\mu) := \min_\pi
G(\pi, \mu)\) is weak-* continuous: by (b) the family
\(\{a \mapsto \mathcal F_a(\pi) : \pi \in \Pi\}\) is uniformly
bounded and uniformly equicontinuous in \(C(\mathscr S)\), so
Arzel\`a--Ascoli makes its closure compact, and \(\mu_n \to \mu\)
weak-* forces \(\sup_\pi |G(\pi, \mu_n) - G(\pi, \mu)| \to 0\);
evaluating at near-minimizing \(\pi\)'s gives \(g(\mu_n) \to
g(\mu)\).

(e) \emph{Atomic approximation.} Finitely supported measures are weak-*
dense in \(\mathcal P(\mathscr S)\): quantize \(\mathscr S\) into
finitely many Borel sets of diameter below \(1/n\), keep one
representative point per set, and give it the set's \(\mu\)-mass; the
resulting finitely supported measures converge to \(\mu\). Taking
\(\mu_n \to \mu^{*}\), the maximizer of (c), (d) gives
\(g(\mu_n) \to g(\mu^{*}) = J_{\mathcal I}(T)\). Hence
\[
J_{\mathcal I}(T) \;=\; \sup\Bigl\{ \min_\pi
\textstyle\sum_\ell \lambda_\ell \mathcal F_{a_\ell}(\pi)
\;:\; q \in \mathbb N,\ a_\ell \in \mathscr S,\
\lambda \in \Delta_q \Bigr\},
\]
which is (9) once the inner minimum is identified with
\(\Gamma_{\mathcal I}\).

(f) \emph{The \(\Gamma\) representation.} For
\(\mu = \sum_\ell \lambda_\ell \delta_{a_\ell}\) the inner minimum
decomposes over information cells \((r, C)\): it is the infimum over
\(u_{r,C} \in L^\infty(I_r; U)\) of
\(\int_{I_r} g_{r,C}(\sigma)^\top u_{r,C}(\sigma)\,d\sigma\) with
\(g_{r,C} = \sum_{\ell:\,\theta_\ell \in C} \lambda_\ell
k_{a_\ell}\). The inequality \(\le \int \phi_U(g_{r,C})\) is the
pointwise bound; the converse is measurable selection --- for a.e.\
\(\sigma\) the set of minimizers of \(v \mapsto g_{r,C}(\sigma)^\top
v\) over the polytope \(U\) is a nonempty face, the lexicographically
first minimizing vertex in a fixed vertex order is measurable in
\(\sigma\), and selecting it attains the pointwise minimum a.e.,
so the infimum equals \(\int_{I_r} \phi_U(g_{r,C})\). Summed over
cells, \(\min_\pi \sum_\ell \lambda_\ell
\mathcal F_{a_\ell}(\pi) = \Gamma_{\mathcal I}(\lambda, a)\) as
defined in (1).

(g) \emph{Nonviability equivalence.} If \(J_{\mathcal I}(T) > 0\), (e)
supplies a finitely supported \(\mu\) with
\(g(\mu) > J_{\mathcal I}(T)/2 > 0\): a positive finite atomic
certificate. Conversely
\(\Gamma_{\mathcal I}(\lambda, a) = \min_\pi \sum_\ell
\lambda_\ell \mathcal F_{a_\ell}(\pi) \le \min_\pi \max_a
\mathcal F_a(\pi) = J_{\mathcal I}(T)\) because the maximum bounds the
\(\lambda\)-weighted average. \(\square\)

\emph{Step 6 (sparsity).} Factor the finite row map through its rank,
\(A v = L(Rv)\), \(Rv \in \mathbb R^{r}\); the attainable moment set
\(\mathcal M = \{Rv : v \in U^{Q}\}\) is compact and convex, and each safety
row defines a halfspace in \(\mathbb R^{r}\). If the common feasible set is
empty, Helly's theorem in \(\mathbb R^{r}\) gives an obstruction involving at
most \(r + 1\) safety halfspaces, and finite duality on that subproblem
yields a sparse positive certificate with the same objective: \(r + 1 \le mQ
+ 1\) labels. To preserve a prescribed margin \(\gamma_0 < \rho_{\mathcal
H,G}\), apply the argument to the still-infeasible shifted constraints \(A_a
v \le \beta_a^{+} + \bar e_a + \gamma_0\). The continuous finite-rank case
follows identically on the compact attainable moment body of the lifted
kernels, giving \(r_{\mathcal I} + 1\). \hfill\(\square\)
\rulebar

\textbf{Proof of Proposition 3 (arbitrarily large minimal obstructions).}
With the construction of the main text and \(w_j = \int_{I_j} u\): mode
\(j \le q\) ends the interval at \(x_j = w_j - 1 \ge 0\), so safety of the
blind window requires \(w_j \ge 1\); mode \(q+1\) ends at \(q - \varepsilon -
\sum_j w_j\), so safety requires \(\sum_j w_j \le q - \varepsilon\). Summing
the \(q\) lower bounds against the upper bound gives \(0 \le -\varepsilon\):
a contradiction, so the full belief is nonviable. Removing mode \(q+1\),
\(u \equiv 1\) gives \(w_j = 1\) and \(x_j = 0\): safe. Removing mode \(j \le
q\), \(u = 0\) on \(I_j\) and \(u = 1\) elsewhere gives \(x_{q+1} =
(q-1) - (q - \varepsilon) + 1 = \varepsilon \ge 0\) at the end and safety on
the remaining branch: every proper subbelief is viable. Equal weights
\(\lambda_j = 1/(q{+}1)\) on the \(q+1\) endpoint rows give pooled pre-kernels
\(\sum_j \lambda_j k_j = 0\) on every interval where the modes agree and a
margin of \(\varepsilon/(q{+}1)\) at the horizon. The kernel rank: the
endpoint rows are \(\mathbf 1_{I_j}\)-weighted copies of the single input
direction with one linear relation (their weighted sum), so the rank of the
safety-row matrix is exactly \(q\), and Helly-type sparsity in that rank
requires \(q + 1\) witnesses --- no bound in the input dimension \(m = 1\)
exists. \hfill\(\square\)
\rulebar

\textbf{Proof of the erosion identity.} \(\phi_U\) is concave and
positively homogeneous; for any decomposition,
\(\phi_U(\sum_j g_j) \ge \sum_j \phi_U(g_j)\) because each
\(\phi_U(g_j) = \min_{v \in U} g_j^\top v\) is attained at a common minimizer
for the sum while separate minima only lower the total. Consequently the same
witness family, evaluated on the coarse partition, dominates its evaluation
on the refinement, \(\Gamma_{\mathrm{coarse}} \ge \Gamma_{\mathrm{fine}}\),
and the exact erosion is the displayed integral (12). For
\(U = [-\bar u, \bar u]^m\), \(\phi_U(g) = -\bar u \|g\|_1\), giving the
stated \(\ell^1\) form. The \(3{,}240\)-pair verification for the hexagonal
\(U\) is reported in S2. \hfill\(\square\)
\rulebar

\textbf{Proof of the oracle-recourse certificate (13).} The oracle relaxation
replaces the true post-\(\tau\) policy class by the larger class that may
assign each stored scenario its own measurable control. The certificate (13)
is the atomic margin (1) evaluated in that enlarged class: before \(\tau\)
the pooled kernel term is \(\phi_U(\sum_\ell \lambda_\ell k_\ell)\), after
\(\tau\) each scenario \(\ell\) contributes \(\phi_U(\lambda_\ell k_\ell)\)
with its own margin \(\beta_\ell\) computed from the stored scenario. A
positive margin certifies infeasibility in the enlarged class by
Propositions 1--2; every policy of the true class is feasible for the
enlarged one, so the verdict transfers downward: no actual output-feedback
policy is safe. Soundness is thus immediate; completeness fails because the
oracle class is strictly larger.
\rulebar

\section*{S2. The verification record}

All numbers in the main text are re-derived in exact rational arithmetic
(Python \texttt{fractions}; no floating point in any witness computation;
floats only inside the solver call, whose output is compared against the
exact value only within solver tolerance). Two scripts carry the record, both
at \url{https://github.com/MIKEAA2020/general-sustainability} (folder
\texttt{arena agent 1/paper rewrites/latex}):

\begin{itemize}\tightlist
\item \texttt{paper2c\_lp\_campaign.py} --- the solver-assisted, exactly certified mesh
campaign: builds program (5) at each mesh \(h \in \{0.2, 0.1, 0.05, 0.02,
0.01\}\), verifies the dimension identities (\(mQ{+}1\) variables;
\(Ms(N_t{+}1) + p_U Q\) rows), solves with HiGHS, and certifies optimality of
\(\rho(h) = \max\{0, \tfrac{3}{50} - Th/4\}\) by the exact primal witness
(pre-observation blocks \(0\); post-observation blocks of mode \(j\) at
\(-n_j\); every row checked) and the exact dual witness
(\(\lambda = (\tfrac38, \tfrac5{16}, \tfrac5{16})\),
\(\mu_{j,k} = \lambda_j a_k \eta_j\); stationarity and objective \(-\rho(h)\)
per block). Also re-derives \(\Gamma(\tau) = \tau - \tfrac7{50}\),
\(\tau_{\max} = \tfrac7{50}\), the \(h = \tfrac1{10}\) Farkas contradiction
\(-\tfrac3{100}\), and the pair-viability positions \(1.9752 = \tfrac{12345}{6250}\) and
\(1.9352 = \tfrac{2419}{1250}\), attained by the recorded pairwise
policies (blind-window hold \(u = 0\) resp.\ \(u = n_1\), then the
two-phase brake \(-\tfrac25 n_j\) for \(\tfrac15\) and \(-n_j\) to
rest; Section~6 of the main text). Current run: 9/9 checks pass; solver agreement at most
\(4.2 \times 10^{-17}\) on all five meshes.
\item \texttt{paper2\_nonstandard\_verification.py} --- the 51-check
exact-arithmetic record for the predecessor material: the three-branch
instance values, the \(m+1\) construction for all tested \(q\), the
\(\phi_U\) superadditivity on \(3{,}240\) rational pairs, the error-bound
checks, and the counterexamples. Current run: 51/51 checks pass.
\end{itemize}

The mesh table of the main text is regenerated verbatim by the first script;
the reported iteration counts and wall-clock times are measured values from
the run recorded here (scipy HiGHS, dual simplex), not estimates.

\end{document}
